% Source: https://artofproblemsolving.com/community/c3h2052375_number_theory
% Stable identity: aops-topic-2052375-q01
Let the number be $N=a\cdot10^k+m$, where $a\in\{1,\dots,9\}$ is the leading digit, $k\ge1$, and $0\le m<10^k$ is the number left after erasing $a$. We need $N=57m$, i.e.
\[a\cdot10^k=56m=7\cdot8\cdot m .\]
Since $7\mid a\cdot10^k$ and $\gcd(7,10)=1$, we get $7\mid a$, so $a=7$. Then $m=\dfrac{10^k}{8}$, which is an integer only for $k\ge3$, giving $m=125\cdot10^{k-3}$, which indeed has $k$ digits. So
\[N=7\cdot10^k+125\cdot10^{k-3}=7125\cdot10^{k-3}.\]
Check: $7125=57\cdot125$. So such numbers exist, and they are exactly
\[\boxed{7125,\ 71250,\ 712500,\ \dots\ \ (7125\cdot10^j,\ j\ge0)}.\]
